\documentclass[11pt,a4paper]{article}
\usepackage[a4paper,margin=22mm]{geometry}
\usepackage{amsmath,amssymb,amsthm}
\usepackage{fancyvrb,multicol}
\usepackage[colorlinks=true,linkcolor=black,urlcolor=blue]{hyperref}
\newtheorem*{theorem}{Theorem}
\title{Spreading 50 Points in the Plane}
\author{Limin Ge\thanks{Found and checked by Claude, a model released by Anthropic, running on the author's laptop. The author sent the prompts.}}
\date{October 9, 2026}
\begin{document}
\maketitle
\begin{abstract}
We place 50 points in the plane so that the squared ratio of the largest to the smallest distance between them is at most 47.93329+. Erich Friedman's table of best known configurations lists n up to 31. There is no larger configuration in the table to compare with. The ratio is computed exactly from the published coordinates. We do not claim it is optimal.
\end{abstract}

\section{Result}
\begin{theorem}
There are 50 points in the plane with
\[ \Bigl(\frac{\max d}{\min d}\Bigr)^2 = 47.93329\ldots \]
\end{theorem}
The exact value, a rational number, is given in Appendix~A.


\section{Context}
Erich Friedman's page \emph{Minimizing the Ratio of Maximum to Minimum Distance} lists the best known configurations for $n\le 31$. The value for $n=50$ is not on the page. There is no larger configuration in the table to compare with.


\section{Verification}
The coordinates are decimals with 12 digits after the point, so they are exact rationals; all squared distances, and their largest-to-smallest ratio, are computed with integers. Run
\begin{verbatim}
python friedman-maxmin2d/verify.py friedman-maxmin2d/coords/n50.txt
\end{verbatim}
in \url{https://github.com/Ge-limin/math}.


\section{How it was found}
Random starts of a constrained local optimizer (SLSQP), then basin hopping seeded from the neighbouring sizes (\texttt{friedman-maxmin2d/}).


\appendix
\section{Exact value}
{\small
\begin{verbatim}
numerator:
9586659079901342879593792
denominator:
199999999999616427203441
\end{verbatim}}

\section{Coordinates}
50 points, $x$ and $y$ per line. File: \texttt{friedman-maxmin2d/coords/n50.txt}.
{\tiny
\begin{multicols}{2}
\VerbatimInput{coordinates.txt}
\end{multicols}}
\end{document}
